\documentclass[a4paper,12pt]{article}
\usepackage{inputenc}[utf8]
\usepackage[margin=2cm, headheight=120pt, tmargin=160pt]{geometry}
\usepackage[polish]{babel}

\usepackage{titlesec}
\usepackage{titling}
\usepackage{fancyhdr}
\usepackage{mdframed}
\usepackage{graphicx}
\usepackage{amsmath}
\usepackage{amsfonts}
\usepackage{multicol}
\usepackage{listings}
\usepackage{caption}
\usepackage{float}
\usepackage{pdfpages}
\usepackage[section]{placeins}

\usepackage{multirow}
\usepackage[table,xcdraw]{xcolor}



\usepackage{hyperref}
\hypersetup{
    colorlinks=true,
    linkcolor=blue,
    filecolor=magenta,      
    urlcolor=cyan,
}

\urlstyle{same}
\graphicspath{ {./ilustracje/} }
\title{Wyznaczanie wartości przyspieszenia ziemskiego metodą spadku swobodnego}
\author{Grzegorz Koperwas}
\date{\today}


\usepackage{polski}
\usepackage[polish]{babel} 
\usepackage{indentfirst}
\usepackage{icomma} 

\brokenpenalty=1000
\clubpenalty=1000
\widowpenalty=1000    

\renewcommand\thesection{\arabic{section}.}
\renewcommand\thesubsection{\arabic{section}.\arabic{subsection}.}
\renewcommand\thesubsubsection{\arabic{subsubsection}.}


\newcommand{\stopnie}{\ensuremath{^{\circ}}}

\newcommand{\twierdzonko}[1]{
    \begin{center}
    \begin{mdframed}
    #1
    \end{mdframed}          
    \end{center}
} 

\newcommand{\dwanajeden}[2]{
\ensuremath \left( \begin{array}{c}
    #1\\
    #2
\end{array} \right)
}  

\pagestyle{fancy}
\fancyhead{}
\fancyfoot{}

\rfoot{}
\lfoot{}
\lhead{}
\rhead{}
\chead{
    \includegraphics[width=\linewidth]{banner.png}
    \scriptsize{\textbf{Politechnika Śląska jako Centrum Nowoczesnego
    Kształcenia opartego o badania i innowacje}}

    \scriptsize{\texttt{POWR-03.05.00-00-Z098/17/00foobar}}

    \vspace{1pt}
    \hrule
    
    \begin{flushleft}
        \begin{tabular}{ r l }
            Nazwa projektu PBL: & foobar \\
            Imię i nazwisko studenta: & {{ author_name }} \\
            Wydział / Kierunek / Specjalność: & Matematyka Stosowana / Informatyka /
            brak \\
            Rok akademicki / semestr: & 2021-2022 / 4 \\
            Opiekun główny projektu: & foobar
        \end{tabular}
    \end{flushleft}
}
\cfoot{\scriptsize\bfseries Projekt współfinansowany przez Unię Europejską ze środków Europejskiego Funduszu Społecznego}
\renewcommand{\headrulewidth}{0pt}
\renewcommand{\footrulewidth}{1pt}

\usepackage{array}
\newcolumntype{L}{>{\arraybackslash}m{10cm}}

\begin{document}

{% for month in data %}

\begin{center}
    \textbf{Karta ewidencji czasu pracy studenta za miesiąc {{ month.name }} }
\end{center}
\vspace{-5mm}

\begin{table}[H]
\centering
\scriptsize
\begin{tabular}{|cllc|L|}
\hline
\rowcolor[HTML]{F1F1CE} 
\multicolumn{4}{|c|}{\cellcolor[HTML]{F1F1CE}{\color[HTML]{333333} \textbf{Harmonogram czasu pracy}}} &
    \multicolumn{1}{c|}{\textbf{Wykonanie}} \\ \hline
\rowcolor[HTML]{F1F1CE} 
\multicolumn{1}{|c|}{\cellcolor[HTML]{F1F1CE}} &
  \multicolumn{2}{c|}{\cellcolor[HTML]{F1F1CE}\textbf{\begin{tabular}[c]{@{}c@{}}rozkład czasu\\ pracy\end{tabular}}} &
  \cellcolor[HTML]{F1F1CE} &
  \cellcolor[HTML]{F1F1CE} \\ \cline{2-3}
\multicolumn{1}{|c|}{\multirow{-2}{*}{\cellcolor[HTML]{F1F1CE}\textbf{\begin{tabular}[c]{@{}c@{}}Dzień\\ miesiąca\end{tabular}}}} &
  \multicolumn{1}{c|}{\textbf{od}} &
  \multicolumn{1}{c|}{\textbf{do}} &
  \multirow{-3}{*}{\cellcolor[HTML]{F1F1CE}\textbf{\begin{tabular}[c]{@{}c@{}}Razem\\ liczba\\ godzin\\ pracy\end{tabular}}} &
      \multicolumn{1}{|c|}{\multirow{-2}{*}{\cellcolor[HTML]{F1F1CE}\textbf{Zakres
    czynności}}} \\[1.2em] \hline
{% for commit in month.commits %}
    \multicolumn{1}{|l|}{ {{commit.day}} } &
    \multicolumn{1}{l|}{ {{commit.start}} } &
    \multicolumn{1}{l|}{ {{commit.stop}} } &
    \multicolumn{1}{l|}{ {{commit.hours}} } &
    {{commit.msg}} \\ \hline
{% endfor %}
\end{tabular}
\end{table}

\vfill

\begin{minipage}{4cm}
............................

\scriptsize{(Data i podpis studenta)}
\end{minipage}
\hfill
\begin{minipage}{4cm}
................................

\scriptsize{(podpis opiekuna głównego)}
\end{minipage}

\vspace{5mm}

\clearpage

{% endfor %}

\end{document}
